% Apéndice D: Basys 2 frente a Basys 3. Fuente: docs/basys2_frente_basys3.md
% (cada cifra lleva allí su fichero de origen en docs/ o results/).

\chapter{Las dos placas FPGA: Basys~2 y Basys~3}
\label{app:fpga}

La fuente de entropía se llevó a dos placas \ac{FPGA}. La primera
disponible (14 a 16 de noviembre de 2026) fue una Digilent Basys~2 con una
Spartan-3E XC3S100E (90\,nm); el diseño se había preparado para una
Basys~3 con una Artix-7 XC7A35T (28\,nm), que llegó el 17 de noviembre y
en la que se repitió casi todo. Los resultados de este trabajo son los de
la Basys~3. Los de la Basys~2 se conservan como una pasada preliminar: lo
que enseñaron se aplicó en la Basys~3, y un experimento (el biestable de
muestreo con y sin habilitación de reloj) solo existe allí. En este
apéndice los dos osciladores en anillo se nombran como en el \ac{RTL}: RO1
es el anillo muestreado y RO2 el de muestreo, y el divisor es
$K_D = 2^{n}$ periodos de RO2 (registro \texttt{CFG\_KD} $= n$). $Q$ es la
varianza de la fase relativa entre los anillos acumulada por periodo de
RO2, en ciclos de RO1 al cuadrado. Todos los valores son \meas{} en placa
salvo que se indique otra cosa.

\begin{table}[htbp]
\centering\small
\caption{Las dos plataformas.}
\label{tab:app_fpga_plat}
\begin{tabular}{@{}>{\raggedright\arraybackslash}p{0.22\linewidth}>{\raggedright\arraybackslash}p{0.35\linewidth}>{\raggedright\arraybackslash}p{0.35\linewidth}@{}}
\toprule
 & Basys~2 & Basys~3 \\
\midrule
Dispositivo & XC3S100E, Spartan-3E, 90\,nm, 1\,920 LUT4 & XC7A35T, Artix-7, 28\,nm \\
Qué cabe & solo la fuente de entropía (anillos, \ac{ERO}, tests de salud,
  \acs{I2C}, FIFO del flujo); la criptografía sola pide unas 6\,800 LUT4
  \est{}, 3,5 veces el dispositivo & el motor entero: 4\,011 LUT (19\,\%) y
  4\,561 biestables (11\,\%) \\
Herramientas & ISE 14.7; XST no lee \acs{VHDL}-2008, así que antes GHDL
  pasa el diseño a Verilog & Vivado 2025.2 \\
Reloj del sistema & 50\,MHz & 100\,MHz \\
Etapas de los anillos (RO2/RO1) & 7/9 & 11/13 \\
Frecuencia de RO2/RO1 & 143,8/108,7\,MHz & 207,2/167,2\,MHz \\
Retardo por etapa (LUT y rutado, $1/(2Nf)$) & 0,50-0,51\,ns & 0,22-0,23\,ns \\
Colocación de los anillos & grupos de área en bordes opuestos, y después
  LOC/BEL de cada etapa sacados de la netlist implementada & una región
  por anillo, y una LUT6 por etapa fijada con LOC/BEL (sin ello Vivado
  juntaba inversores de RO2 en una LUT6 y el anillo iba a 130\,MHz) \\
Biestable de muestreo & con y sin habilitación de reloj & sin habilitación de reloj \\
Cortes de los tests de salud & 22/596, después 39/683 y 51/766 & 51/766 \\
Divisor mínimo (\texttt{G\_KD\_MIN}) & 7 & 19 \\
\bottomrule
\end{tabular}
\end{table}

\begin{table}[htbp]
\centering\small
\caption{Resultados lado a lado. ``Silenciosa'' es el envoltorio sin
display ni frecuencímetro, el más parecido al circuito integrado; en la
Basys~3, ``libre'' es sin fijar la colocación de las etapas y ``fija'',
fijándola.}
\label{tab:app_fpga_res}
\begin{tabular}{@{}>{\raggedright\arraybackslash}p{0.22\linewidth}>{\raggedright\arraybackslash}p{0.35\linewidth}>{\raggedright\arraybackslash}p{0.35\linewidth}@{}}
\toprule
 & Basys~2 & Basys~3 \\
\midrule
$Q$, ajuste de la autocorrelación $c(M)$ & $1{,}0$--$4{,}9\times10^{-5}$
  (silenciosa, $n = 7$--$9$, según la implementación; tres capturas
  dentro de $\pm3$--$5$\,\%) & $1{,}42\times10^{-5}$ (silenciosa, fija,
  $n = 2$); $0{,}32$--$2{,}75\times10^{-5}$ entre variantes (tres capturas
  dentro de $\pm0{,}3$\,\%) \\
Suelo de ruido blanco de $Q$ (espectro de la frecuencia instantánea) &
  no medido (sin buffer de captura, sin $n \le 2$) &
  $0{,}6$--$1{,}1\times10^{-6}$ ($\sigma \approx 4$--$5$\,ps por periodo
  de 4,9\,ns; calibración \simu{}) \\
Lógica colgada de los anillos & cambia $Q$ $\times2{,}5$ & cambia $Q$
  $\times3$--$4{,}4$, también con las etapas fijadas \\
SP~800-90B, $10^6$ bits consecutivos & $n = 12$: 0,760; $n = 13$--$16$:
  0,82-0,86 & $n = 12$--$19$: 0,815-0,879, saturada; bitstream final a
  $n = 19$: 0,853 \\
SP~800-90B, 8\,378\,368 bits & 0,890 ($n = 14$) & 0,932 (libre, $n = 16$),
  0,880 (fija, $n = 16$) \\
Fracción de unos & 0,496-0,498 (silenciosa), 0,519 (con display) &
  0,495-0,497 (libre y final), 0,505-0,506 (fija), 0,513-0,519 (con
  display): el signo cambia con la implementación \\
Correlación a retardo 1, $\rho_1$ & 0,003-0,006 ($z$ hasta 13) con
  habilitación; $\le 0{,}0005$ ($z \le 1$) sin ella & $|z| \le 2{,}1$ (sin
  habilitación) \\
Tests de salud en placa & test de repetición a $n = 8$, 11 y 12, y de
  proporción a $n = 8$ (22/596); solo el de repetición a $n = 10$
  (51/766) & test de repetición a $n = 8$ (libre); ninguno ni a $n = 8$
  (fija, donde SP~800-90B da 0,186) \\
Prueba de reinicios (SP~800-90B \S3.1.4, 1\,000 reconfiguraciones) & pasa
  a $n = 14$ ($H_r$ 0,855, $H_c$ 0,873, $H_I$ 0,845) & pasa a $n = 16$
  (0,843, 0,892, 0,819) y a $n = 19$ (0,854, 0,830, 0,857) \\
Motor entero con el ESP32 de oráculo & no cabe & 40/40 etiquetas
  \acs{CMAC} iguales a las de mbedTLS, en tres sesiones \\
Divisor de trabajo & $n \ge 16$, provisional (del ajuste de $c(M)$) &
  $n = 19$ (del suelo de ruido blanco) \\
\bottomrule
\end{tabular}
\end{table}

\paragraph{Lo común.} En dos familias y dos nodos el cuadro es el mismo. El
ajuste de $c(M)$ da un ruido de fase relativo de $10^{-6}$--$10^{-5}$ por
periodo de RO2, y el modelo del \ac{ERO} deja residuos de varias a decenas
de veces el ruido binomial: parte de la fase es una modulación
determinista, no jitter. $Q$ depende de la lógica colgada de los anillos y
de la implementación, en el mismo sentido en las dos placas (sin el
frecuencímetro $Q$ es mayor). SP~800-90B se satura en 0,82-0,88 con
$10^6$ bits y no separa la \ac{FPGA} del ESP32 ni de bits ideales, ni, en
la Basys~3, un divisor con $H_\infty = 0{,}47$ de otro con 0,9996. El
sesgo es del 0,2-1,9\,\% y su signo cambia entre implementaciones en las
dos placas. Los tests de salud no ven el régimen de poco jitter, así que
solo lo protege el divisor mínimo. La prueba de reinicios pasa en las dos.

\paragraph{Qué cambió el divisor.} No la placa, sino el método. El ajuste
de $c(M)$ cuenta la modulación como jitter y sobreestima $Q$; el método
espectral, posible solo en la Basys~3 porque necesita el buffer de captura
a divisores pequeños, separa el suelo blanco, entre 13 y 24 veces por
debajo del ajuste en la misma variante. Con él, el divisor de la \ac{FPGA}
pasó de $n \ge 16$ a $n = 19$. La cota $n \ge 16$ de la Basys~2 no se
puede confirmar ni descartar. El mismo $n$ no es el mismo tiempo en las
dos placas: $n = 16$ a 143,8\,MHz son 456\,$\mu$s por bit, y $n = 19$
a 205\,MHz, 2,56\,ms.

\paragraph{Lo que solo muestra la Basys~2.} Solo allí se comparó el
biestable de muestreo con y sin habilitación de reloj, con los anillos y
el biestable en los mismos sitios: con ella, $\rho_1 \approx 0{,}005$; sin
ella, nada. Por eso el motor muestrea sin habilitación por defecto
(Capítulo~\ref{ch:ch4_engine}). La Basys~3 no tiene variante con
habilitación, así que la causa descansa en una sola familia; la ausencia
de correlación sin ella se confirma en las dos.

\paragraph{Lecciones aplicadas.} La Basys~2 mostró que un anillo puede
desaparecer en la implementación (el mapeador colapsó la cadena de
inversores), que los anillos colocados van más deprisa de lo que cierra el
contador del \ac{ERO}, que el display multiplexado modula los anillos, y
que la repetibilidad hay que medirla antes de leer una diferencia. En la
Basys~3 el script de implementación lista cada celda de los anillos tras
rutar y se para si falta un anillo, el número de etapas y la restricción de
tiempos de RO2 siguen a la frecuencia medida, las campañas usan el
envoltorio silencioso y cada $Q$ sale de tres capturas.

\paragraph{90\,nm frente a 28\,nm.} Colocados, los anillos de la Artix-7
son unas 2-2,3 veces más rápidos por etapa. El ruido de fase es del mismo
orden en las dos placas, pero con frecuencias, divisores, alcances del
ajuste y envoltorios distintos, y con un modelo que no describe su forma en
ninguna; no se concluye qué nodo tiene más jitter. La única cifra que
aspira a ser térmica, el suelo blanco de la Basys~3 ($\sigma \approx
4$--$5$\,ps por periodo de 4,9\,ns), es del orden de los 2-4\,ps por
periodo que se publican para osciladores en anillo en
28\,nm~\cite{petura_2016_survey}.
